\section*{अध्याय \ref{ch:b1:crystals-metals} --- क्रिस्टल I: आदर्श क्रिस्टल और धातुएँ}

\begin{solution}{exo:b1:crystals-metals:1}
सरल घनीय: $8 \times \frac18 = 1$ परमाणु, उपसहसंयोजन 6। BCC: $8 \times
\frac18 + 1 = 2$, उपसहसंयोजन 8। FCC: $8 \times \frac18 + 6 \times
\frac12 = 4$, उपसहसंयोजन 12।
\end{solution}

\begin{solution}{exo:b1:crystals-metals:2}
परमाणु कोर के अनुदिश छूते हैं: $a = 2R$, $Z = 1$, $C = \frac43\pi R^3/(2R)^3
= \pi/6 \approx 0.52$, जो BCC (0.68) और FCC (0.74) से कम है।
\end{solution}

\begin{solution}{exo:b1:crystals-metals:3}
$R = a\sqrt2/4 = 405.0 \times 0.3536 = \qty{143.2}{pm}$;
$\rho = 4 \times 26.98/(6.022 \times 10^{23} \times (4.050 \times
10^{-8})^3) = \qty{2.70}{g/cm^3}$।
\end{solution}

\begin{solution}{exo:b1:crystals-metals:4}
$R = a\sqrt3/4 = \qty{185.8}{pm}$; $\rho = 2 \times 22.99/(6.022 \times
10^{23} \times (4.291 \times 10^{-8})^3) = \qty{0.966}{g/cm^3}$, जो
मापे गए \qty{0.97}{g/cm^3} से मेल खाता है।
\end{solution}

\begin{solution}{exo:b1:crystals-metals:5}
$a^3 = 4M/(\rho N_A) = 4 \times 107.87/(10.50 \times 6.022 \times 10^{23}) =
\qty{6.824e-23}{cm^3}$, इसलिए $a = \qty{4.086e-8}{cm} = \qty{408.6}{pm}$ और
$R = a\sqrt2/4 = \qty{144.5}{pm}$।
\end{solution}

\begin{solution}{exo:b1:crystals-metals:6}
उपपत्ति \cref{prop:b1:crystals-metals:hcp} वाली ही है। मैग्नीशियम:
$c/a = 521.0/320.9 = 1.624$। आयतन
$\frac{3\sqrt3}{2}a^2c$ वाले प्रिज़्म में 6 परमाणु हैं, इसलिए $\rho = 6M/(N_A \cdot
\frac{3\sqrt3}{2}a^2c) = 6 \times 24.31/(6.022 \times 10^{23} \times
2.598 \times (3.209 \times 10^{-8})^2 \times 5.210 \times 10^{-8}) =
\qty{1.74}{g/cm^3}$।
\end{solution}

\begin{solution}{exo:b1:crystals-metals:7}
परमाणु: $(0,0,0)$, $(\frac a2,\frac a2,0)$, $(\frac a2,0,\frac a2)$,
$(0,\frac a2,\frac a2)$। \omterm{def:b1:crystals-metals:site}{अष्टफलकीय स्थल}: $(\frac a2,\frac a2,\frac a2)$,
$(\frac a2,0,0)$, $(0,\frac a2,0)$, $(0,0,\frac a2)$। \omterm{def:b1:crystals-metals:site}{चतुष्फलकीय स्थल}: वे
8 बिंदु जिनका हर निर्देशांक $\frac a4$ या $\frac{3a}{4}$ है। प्रति परमाणु एक
अष्टफलकीय और दो \omterm{def:b1:crystals-metals:site}{चतुष्फलकीय स्थल}।
\end{solution}

\begin{solution}{exo:b1:crystals-metals:8}
$R(\ce{Cu}) = 361.5\sqrt2/4 = \qty{127.8}{pm}$, $R(\ce{Ni}) = 352.4\sqrt2/4
= \qty{124.6}{pm}$: त्रिज्याएँ 3\,\% के भीतर हैं और संरचना एक ही है, इसलिए कोई भी
परमाणु \omterm{def:b1:crystals-metals:lattice}{जालक} पर दूसरे का स्थान ले सकता है। हाइड्रोजन, जो कहीं छोटा है,
केवल \omterm{def:b1:crystals-metals:site}{अंतराकाशी स्थलों} में बैठ सकता है: वह \omterm{def:b1:crystals-metals:alloy}{अंतराकाशी मिश्रधातु} बनाता है।
\end{solution}

\begin{solution}{exo:b1:crystals-metals:9}
$R = 316.5\sqrt3/4 = \qty{137.0}{pm}$; $\rho = 2 \times 183.84/(6.022 \times
10^{23} \times (3.165 \times 10^{-8})^3) = \qty{19.26}{g/cm^3}$; प्रति
\unit{cm^3} परमाणु: $2/a^3 = 2/(3.165 \times 10^{-8})^3 = \num{6.31e22}$।
\end{solution}

\begin{solution}{exo:b1:crystals-metals:10}
फलक-केंद्र $(\frac a2,\frac a2,0)$ ऊपर और नीचे की कोष्ठिकाओं के अंतःकेंद्रों से
$a/2$ दूर है: $r_O = \frac a2 - R = \frac a2 - \frac{a\sqrt3}{4}
= \frac{2-\sqrt3}{4}a$। बिंदु $(\frac a2,\frac a4,0)$ की दूरी
$\sqrt{a^2/4 + a^2/16} = \frac{a\sqrt5}{4}$ है, कोनों $(0,0,0)$ और
$(a,0,0)$ से भी और अंतःकेंद्रों $(\frac a2,\frac a2,\pm\frac a2)$ से भी:
$r_T = \frac{\sqrt5 - \sqrt3}{4}a$। $a = 4R/\sqrt3$ लेने पर:
$r_O = \frac{2-\sqrt3}{\sqrt3}R = 0.155\,R$ और
$r_T = \frac{\sqrt5-\sqrt3}{\sqrt3}R = 0.291\,R$। BCC में \omterm{def:b1:crystals-metals:site}{चतुष्फलकीय स्थल}
बड़ा है।
\end{solution}

\begin{solution}{exo:b1:crystals-metals:11}
$C = Z\cdot\frac43\pi R^3/a^3$ और $Z/a^3 = \rho N_A/M$, इसलिए $C =
\frac{4\pi R^3\rho N_A}{3M} = \frac{\pi\rho N_A}{6M}(2R)^3$। कॉपर:
$\pi \times 8.93 \times 6.022 \times 10^{23}/(6 \times 63.55) \times (2.556
\times 10^{-8})^3 = 0.740$।
\end{solution}

\begin{solution}{exo:b1:crystals-metals:12}
$a \approx (407.8 + 408.6)/2 = \qty{408.2}{pm}$; माध्य मोलर द्रव्यमान $(196.97 +
107.87)/2 = \qty{152.42}{g/mol}$; $\rho = 4 \times 152.42/(6.022 \times
10^{23} \times (4.082 \times 10^{-8})^3) = \qty{14.9}{g/cm^3}$।
\end{solution}

\begin{solution}{pb:b1:crystals-metals:1}
\textbf{1.} $Z = 2$; उपसहसंयोजन 8।
\textbf{2.} मुख्य विकर्ण के अनुदिश: $a\sqrt3 = 4R$, $R = 286.65 \times
\sqrt3/4 = \qty{124.1}{pm}$।
\textbf{3.} $\rho = 2 \times 55.845/(6.022 \times 10^{23} \times (2.8665
\times 10^{-8})^3) = \qty{7.87}{g/cm^3}$।
\textbf{4.} $2/a^3 = 2/(2.8665 \times 10^{-8})^3 = \num{8.49e22}$ परमाणु।
\textbf{5.} $C = \pi\sqrt3/8 = 0.680$।
\textbf{6.} परिकलित और मापे गए घनत्व तीन सार्थक अंकों तक मेल खाते हैं।
\textbf{7.} $\alpha$: $R = 289.8\sqrt3/4 = \qty{125.5}{pm}$; $\gamma$:
$R = 363.9\sqrt2/4 = \qty{128.7}{pm}$।
\textbf{8.} $\alpha$: प्रति परमाणु $a^3/2 = \qty{12.17e6}{pm^3}$; $\gamma$:
प्रति परमाणु $a^3/4 = \qty{12.05e6}{pm^3}$।
\textbf{9.} $\gamma$ अधिक घना है: $\rho_\alpha = 55.845/(6.022 \times 10^{23}
\times 1.217 \times 10^{-23}) = \qty{7.62}{g/cm^3}$, $\rho_\gamma =
\qty{7.70}{g/cm^3}$।
\textbf{10.} $(12.05 - 12.17)/12.17 = -1.0\,\%$: संक्रमण के पार गर्म करने पर
छड़ लगभग 1\,\% सिकुड़ती है।
\textbf{11.} हाँ: FCC अधिक संहत है (0.74 बनाम 0.68), जो बड़ी
परमाणु त्रिज्या पर भारी पड़ता है: $(0.680/0.740) \times (128.7/125.5)^3 = 0.99$।
\textbf{12.} $r_C = a\sqrt3/8 = 356.68 \times 0.2165 = \qty{77.2}{pm}$।
\textbf{13.} $r_O = 0.0670 \times 289.8 = \qty{19.4}{pm}$।
\textbf{14.} $r_T = 0.1260 \times 289.8 = \qty{36.5}{pm}$।
\textbf{15.} कार्बन (\qty{77}{pm}) सबसे बड़े स्थल से दोगुना बड़ा है:
पड़ोसी आयरन परमाणुओं को धकेलकर दूर करना पड़ता है।
\textbf{16.} इस विकृति में बहुत ऊर्जा लगती है, इसलिए बहुत कम कार्बन परमाणु
$\alpha$-आयरन में प्रवेश करते हैं।
\textbf{17.} 4 परमाणुओं वाली कोष्ठिका में 4 \omterm{def:b1:crystals-metals:site}{अष्टफलकीय स्थल}: प्रति आयरन परमाणु एक।
\textbf{18.} $r_T = (\sqrt{3/2} - 1) \times 128.7 = \qty{28.9}{pm}$।
\textbf{19.} प्रति आयरन एक कार्बन: \ce{FeC}, $w(\ce{C}) = 12.011/(55.845 +
12.011) = 17.7\,\%$।
\textbf{20.} दस आयरन पर एक कार्बन: $w = 12.011/(558.45 + 12.011) =
2.1\,\%$।
\textbf{21.} $\gamma$-आयरन के \omterm{def:b1:crystals-metals:site}{अष्टफलकीय स्थल} $\alpha$-आयरन के किसी भी स्थल से
कहीं बड़े हैं: गर्म आयरन कार्बन घोल लेता है, जिसे शमन फिर
फँसा लेता है।
\textbf{22.} $r_O = (\sqrt2 - 1) \times 128.7 = \textbf{\qty{53}{pm}}$, जो
आयरन का सबसे बड़ा रिक्त स्थान है और फिर भी कार्बन परमाणु (\qty{77}{pm}) से छोटा है:
कार्बन गर्म आयरन में घुलता है, पर उसे फैला देता है।
\end{solution}
